%======================================================================
\subsection{BW-SIM-04 --- Built, 1 Nov: ccx error text, folded Tet10 elements, elements outside every solid}\label{sec:sim-ccx-diag-built}
%======================================================================

Three reports from the VL-P01 turbojet project, each a way a study died with no usable message.

\paragraph{1. Empty error text.} CalculiX prints its \code{*ERROR} lines to \emph{stdout}, wrapped over two lines and once per
integration point; stderr stays empty, so \code{CalculiX finished with exit code 201:} ended in nothing. The new
\code{bruceworx\_sim.ccx\_error\_lines} reads the run's output (and, when stdout has none, the \file{.dat}, \file{.sta} and
\file{.cvg} files beside the input), joins each wrapped message, drops repeats, merges the
\code{nonpositive jacobian determinant in element N} lines into one line with the element numbers (eight shown, then the total)
and keeps at most four lines of 200 characters. \code{ccx\_failure\_text} builds the message; the blocking \code{sim.solve}, the
blocking and the QProcess path of the Study panel (the panel now keeps the process output) raise or show it, and the Log pane gets
it as an error message. With no \code{*ERROR} line the text says so and points to the \file{.dat} and \file{.sta} files.

\paragraph{2. Folded Tet10 elements.} Netgen's curved second-order mid-edge nodes fold inside small fillets (R1.5 at
MaxSize 1.2 on the impeller). \code{tet10\_jacobian\_ratios} computes, vectorised in chunks of 60\,000 elements, the
Jacobian determinant of every Tet10 at 15 points (4 corners, 6 mid-edges, the 4 Gauss points of ccx's C3D10, the centroid) and
divides the smallest by the determinant of the straight-sided tetrahedron: 1 for a straight element, at most 0 where ccx stops.
The quotient does not depend on the node order's sign (FemMesh and the CalculiX input differ in it). Elements at or under 0.1
get their six mid-edge nodes moved to the edge midpoints (\code{straighten\_tet10}); a mid node is one node of all elements
sharing the edge, so neighbours stay consistent, and the check repeats (up to four passes) because a neighbour can get worse.
\begin{itemize}
  \item After meshing, \code{\_mesh\_study} checks the FemMesh (\code{repair\_fem\_tet10}) and rebuilds it with the moved
        nodes only when some are bad; the count goes to \code{stats['notes']}, the progress text and the Log.
  \item After the input is written, \code{prepare\_solve} checks the deck itself (\code{repair\_inp\_tet10}, C3D10 elements, any
        node/element order), so a mesh not made by the study (loaded, imported) is covered; folded nodes are rewritten in the file.
  \item If straightening leaves a non-positive Jacobian (flat elements), the study stops with an error naming the elements.
\end{itemize}
\textbf{Repro on the user's deck} (\file{cad/sim/work/Mesh.inp}, 268\,153 Tet10, 439\,000 lines): ccx stopped at 14 elements
(176165, 186887, 219565, 228381, 267476, 281895, 287497, 288966, 298208, 356588, \ldots). The check finds \textbf{570 with a
non-positive Jacobian and 1040 at or under 0.1} (worst $-1.75$), among them every element ccx named; the 5\,875 mid-edge nodes
of those 1040 are straightened; afterwards 0 are at or under 0.1 (worst 0.10) and ccx runs the same deck past the Jacobian
stage into the solver iterations (it was stopped there; the iterative-Cholesky run takes long). The check takes 5.3\,s for 268\,000
elements on this machine, including reading the file.

\paragraph{3. No material assigned.} FEM gives an element to a material solid only if \emph{all} its nodes belong to that
solid, and \code{FemMesh.getNodesBySolid} takes a node as belonging when its distance to the solid is below the solid's
tolerance (about $10^{-7}$). Curved second-order mid-edge nodes at the M8 thread (rev J of the shaft) lie off the surface by
more, so the five elements (67731, 78117, 79943, 97143, 97144 in \file{cad/modal/Mesh.inp}, all sharing two nodes near
$(1.0\ldots1.5,\,-4.2,\,15.9\ldots16.0)$) were in no solid's set and got no material. That the earlier run solved is the old
shaft front having no thread there. The fix is in the FEM core, small and better for everyone:
\code{meshtools.assign\_unreferenced\_femelements}, called at the end of \code{get\_femelement\_sets} when no object takes the
rest: every volume element no reference solid holds goes to the object whose solid is nearest to its centroid (0 when inside).
It warns in the report view and keeps the result in \code{meshtools.LAST\_UNREFERENCED}; \code{prepare\_solve} reports the count in
the Log. A script that uses \code{FemToolsCcx} directly (like \file{modal\_rotor.py}) gets the fix too.

\paragraph{Tests.} \code{TestBruceWorxSimCore.TestCcxDiagnostics}: the error-line merge and trim, the \file{.dat} fallback, the
no-error text, a straight Tet10 (ratio 1), one with a mid node pushed through (ratio $\le 0$, straightened to the midpoint),
a shared mid node moved once for both elements, a flat element reported as remaining, a written deck repaired, and a real
mesh bent and repaired. \code{TestBruceWorxSimRotor.TestMaterialSetsOutsideSolids}: two boxes as a CompSolid, one mid node of
the outer face pushed 0.5\,mm out, every element ends in exactly one material set (11 went to the nearest solid).

\paragraph{Open.} The threshold 0.1 straightens mildly curved elements as well as folded ones (1040 of 268\,000 here); the surface
loses at most the mid-node bulge of a few elements. A proper cure is a finer mesh at the fillet (size field) or Gmsh with
\code{Mesh.SecondOrderLinear}, which has no binary yet. The nonlinear buckling run still shows only the tail of its output.
